\documentclass[12pt]{article}
\usepackage{graphicx} % Required for inserting images
\usepackage[margin=0.5in]{geometry}
\usepackage{amsmath, amssymb}
\usepackage{mathrsfs}


\usepackage{soul}


% Dr. David Brown Commands
\newcommand{\ds}{\displaystyle}
\newcommand{\R}{\mathbb{R}}
\newcommand{\N}{\mathbb{N}}
\newcommand{\Q}{\mathbb{Q}}
\newcommand{\C}{\mathbb{C}}


% Commands to get script r
\usepackage{calligra}
\DeclareMathAlphabet{\mathcalligra}{T1}{calligra}{m}{n}
\DeclareFontShape{T1}{calligra}{m}{n}{<->s*[2.2]callig15}{}
\newcommand{\scripty}[1]{\ensuremath{\mathcalligra{#1}}}

% Unit commands
\newcommand{\degree}{\text{°}}
\newcommand{\unitSpeed}{\frac{\text{m}}{\text{s}}}
\newcommand{\unitAcceleration}{\frac{\text{m}}{\text{s}^2}}

% Constant commands
\newcommand{\avogadro}{6.022\times10^{23}}

\newcommand{\boltzmannConstK}{1.381\times10^{-23}\frac{\text{J}}{\text{K}}}
\newcommand{\boltzmannConstInEV}{8.617\times10^{-5}\frac{\text{eV}}{\text{K}}}
\newcommand{\planckConst}{6.626\times10^{-34}\text{J}\cdot\text{s}}

\newcommand{\atmP}{1.01325\times10^5\,\text{Pa}}

% Known Value Commands
\newcommand{\electronMass}{9.109\times10^{-31}\,\text{kg}}
\newcommand{\protonMass}{1.673\times10^{-27}\,\text{kg}}

% Commands to easily write partial derivatives
\newcommand{\partDeriv}[2]{\frac{\partial #1}{\partial #2}}
\newcommand{\doublePartDeriv}[2]{\frac{\partial^2 #1}{\partial^2 #2}}


\title{Problem 7.52}
\author{Gabriel Decker}


\begin{document}


\maketitle
\begin{flushleft}



In this problem, we will consider a the human body to be a blackbody and do some comparisons with other objects.
We will just need the blackbody intensity equation.
\begin{equation}
    P/A=\sigma eT^4
\end{equation}
with $\sigma=5.67\times10^{-8}\,\text{W/m}^2\text{K}^4$.\\~\\~\\


\textbf{Part a}\\
First, we'll find the total power radiating through our bodies.
This is the following.
$$P=\sigma eAT^4$$
The normal human body temperature is $T=310.15\,\text{K}$.
The average human man surface area is $A=1.9\,\text{m}^2$.
At this low of temperature, the blackbody spectrum is dominated by thermal radiation, which our bodies radiate like a blackbody., so $\epsilon=1$.
So, this gives the following power.
$$\implies P=(5.67\times10^{-8}\,\text{W/m}^2\text{K}^4)(1.9\,\text{m}^2)(310.15\,\text{K})^4$$
$$\implies P=996.8\,\text{W}$$\\~\\~\\



\textbf{Part b}\\
Now considering the daily energy input and output, based on the above power, we can multiply it by one day to get the total energy output $U$ of a person in a vacuum.
$$U=P\cdot86400\,\text{s}$$
$$\implies U=996.8\,\text{J/s}\cdot86400\,\text{s}$$
$$\implies U=8.6124\times10^7\,\text{J}$$\\~\\

The average man should consume about 2500 Calories (or 2500 kcal) of energy from food.
The conversion factor from kcal to joules is $1\,\text{kcal}=4184\,\text{J}$.
This implies that we consume about $1.0460\times10^7\,\text{J}$ of energy from food.
This is about 12\% of the energy that we would need to stay at a constant temperature.
In this problem, we have been assuming that we aren't absorbing anything from our environment.
Obviously, we live in environments that also radiate energy, so we get lots of energy from that.\\~\\~\\



\textbf{Part c}\\
Let's now compare the power per unit mass of humans vs the sun.
We're given that for the sun, $M=2\times10^{30}\,\text{kg}$ and $P=3.9\times10^{26}\,\text{W}$.
For the average human man, we have the above $P=996.8\,\text{W}$ and the looked up $78.1\,\text{kg}$.
With these values, we can compare the values for us against the sun.\\~\\

For the sun, $P/M=0.000195\,\text{W/kg}$.
For the the average man, $P/M=12.8\,\text{W/kg}$.






\end{flushleft}

\end{document}
